\documentclass[twocolumn]{article}
\newcommand{\filler}{Filler text expands from a macro, so every use of the macro is attributed to the line where it is used, not where it is defined. This makes one source line produce several lines of output in a column.}
\begin{document}
\section{Two Columns}
Two-column layout puts lines of text side by side,
so a point in the right column must never be attributed to the left column.
Each source line here is a sentence of moderate length.
The paragraphs are long enough to fill both columns of the first page
and to continue on a second page.

Typesetting systems break paragraphs into lines by minimising a measure of badness
over the whole paragraph rather than line by line.
The result is that a change near the end of a paragraph may move words near its beginning.
Synchronisation therefore records, for every line of output, the source line that produced each word.

\filler

\section{Boxes and Glue}
Every line of output is a horizontal box.
Between the words there is glue that can stretch or shrink.
Kerns adjust the space between particular letter pairs, as in AV or To.
Mathematics such as $x_1 + x_2 = y$ adds math nodes around the formula.
Vertical boxes stack the lines of a paragraph and the paragraphs of a page.

\filler
\filler

\section{Columns}
The left column is filled first; when it is full, the output routine starts the right column.
After both columns are full, the page is shipped out and a new page begins.
Footnotes and floats complicate this picture,
but this document keeps to plain paragraphs and section headings.

A paragraph in the middle of the text with a short line.

Another paragraph that talks about nothing in particular, but it does so at some length,
in order to occupy enough space that the column break falls somewhere inside a paragraph
rather than neatly between two of them, which is the more interesting case to test.
The words continue for a while longer.
They keep going until the column is very likely to have been broken in the middle of this paragraph.

\section{More Text}
\filler

Short.

Medium sized paragraph with a handful of words in it, written on one line.

\filler
\filler
\filler

\section{Second Page}
\filler
\filler \filler

\filler
\filler

A paragraph between the fillers, with words of its own.
It is two source lines long.

\filler
\filler
\filler
\filler

\section{End}
The last section of the two-column document.
It ends here.
\end{document}
